\scanpage{110}

\[
T\longrightarrow Y
\]
\uline{et} d'un relèvement en $T\longrightarrow X$, ou ce qui revient au $=$, d'une section de $X_T=X\times_YT$ sur $T$. Ceci posé;

\diagram{\begin{tikzcd}
X\arrow[d] & X_T\arrow[l]\arrow[d] & \\
Y & T\arrow[l] & T\arrow[ul]\arrow[l,equal]
\end{tikzcd}}

on aura un isom \unsure{can|can\textsuperscript{t}} (sous l'hyp. de ``suite exacte d'homotopie'' faite)
\[
(12)\qquad \mathcal E_{X/Y}^T\simeq\pi_1(X_T;T)\simeq\mathcal E_{X_T}^T
\]
et on aura
\[
(13)\qquad 1\longrightarrow\mathcal E_{X/Y}^T\unsure{\mathrel{\text{---}}|\longrightarrow}\mathcal E_X^T\unsure{\mathrel{\text{---}}|\longrightarrow}\mathcal E_Y^T\unsure{\mathrel{\text{---}}|\longrightarrow}1
\]
(Cette hypothèse sera tjrs satisfaite si $Y=\operatorname{Spec}K$, $K$ un corps, si $X$ est géom. $0$-connexe sur $K$)

Plus généralement, si on a une suite exacte d'homotopie universelle, mais pas néc. avec $\mathcal E_X^T\unsure{\mathrel{\text{---}}|\longrightarrow}\mathcal E_Y^T$ surj,

\margin{109}
